%
% Project: Ethical Alignment of Small Language Models — M1 Advanced Deep Learning
% Florian NGO — Télécom Saint-Étienne
%
% This file is meant to be pasted into the ACL Overleaf template:
% https://www.overleaf.com/latex/templates/association-for-computational-linguistics-acl-conference/jvxskxpnznfj
% Replace the template's main.tex content with this.
%

\title{Comparing Direct, Reinforcement, and Retrieval-Based Alignment of a Small Language Model on the ETHICS Benchmark}

\author{Florian NGO \\
  Télécom Saint-Étienne, IMT Mines-Télécom \\
  \texttt{ngoflorian@gmail.com}}

\begin{document}
\maketitle

\begin{abstract}
We compare three families of ethical alignment methods—Direct Preference Optimization (DPO), Reinforcement Learning from Human Feedback (RLHF), and Retrieval-Augmented Generation (RAG)—applied to a 1.5B-parameter language model under tight compute constraints. All three methods are trained or constructed using freely available preference data (PKU-SafeRLHF and UltraFeedback) and evaluated on the ETHICS benchmark \citep{hendrycks2021ethics}, which is held out entirely from training. Our central finding is that retrieval-augmented alignment, despite involving no parameter updates, outperforms both DPO and RLHF on overall ETHICS accuracy (+9.6 vs +6.0 vs +4.2 points over the baseline). We further observe a plateau effect when retrieval is composed with preference-based fine-tuning, suggesting that the two alignment paradigms encode largely overlapping signals on this benchmark. We discuss implications for compute-constrained alignment of small models.
\end{abstract}

\section{Introduction}

Aligning language models to behave consistently with human ethical values is now a central concern in NLP research and deployment \citep{ouyang2022instructgpt,bai2022constitutional}. While most contemporary work targets models in the 7B--70B parameter range, smaller models (1--3B) remain practically important: they are deployable on consumer hardware, can be fine-tuned without industrial compute, and serve as a stress test for alignment methods that must operate with limited capacity.

This paper reports an empirical comparison of three alignment paradigms applied to a small (1.5B) language model:
\textbf{Direct Preference Optimization (DPO)} \citep{rafailov2023dpo}, which optimizes a policy directly from preference pairs;
\textbf{Reinforcement Learning from Human Feedback (RLHF)} \citep{ouyang2022instructgpt}, comprising reward model training followed by policy optimization via Proximal Policy Optimization (PPO);
and \textbf{Retrieval-Augmented Generation (RAG)} \citep{lewis2020rag}, which injects retrieved ethical principles into the prompt at inference time without any parameter updates.

The methods occupy different points in the alignment design space: DPO and RLHF modify model weights using preference signals; RAG provides external context without retraining. We hypothesize that on a small model with limited expressive capacity, external retrieval may provide a stronger signal than fine-tuning on noisy preference data, particularly for tasks requiring explicit ethical reasoning.

We make the following contributions:
\begin{enumerate}
    \item We implement and compare three alignment paradigms on \texttt{Qwen2.5-1.5B-Instruct} using a unified evaluation protocol on the ETHICS benchmark across five moral foundations.
    \item We report that retrieval-augmented alignment outperforms both preference-based methods overall, achieving +9.6 absolute points over the baseline, while DPO and RLHF yield +6.0 and +4.2 respectively.
    \item We document a \emph{plateau effect}: composing retrieval with preference fine-tuning (DPO+RAG, RLHF+RAG) does not improve over RAG alone, suggesting redundancy between the two alignment signals on ETHICS.
    \item We provide open-source code and ablations\footnote{\url{https://github.com/FeelTheFloww/adl-ethics}}.
\end{enumerate}

\section{Related Work}

\paragraph{Preference-based fine-tuning.} RLHF, as popularized by \citet{ouyang2022instructgpt}, trains a reward model on human preferences and then uses PPO \citep{schulman2017ppo} to optimize a policy against this reward while constraining KL divergence from a reference policy. DPO \citep{rafailov2023dpo} shows that this two-stage pipeline can be replaced by a single supervised loss derived from the same preference data, eliminating the reward model and the PPO sampling loop. Subsequent work \citep{azar2024ipo,xu2024dpo_failures} has examined the relative strengths and failure modes of the two approaches, particularly at small model scale.

\paragraph{Retrieval-augmented generation.} RAG \citep{lewis2020rag} augments a language model's context with documents retrieved from an external corpus. While originally proposed for knowledge-intensive tasks, RAG has been explored as a means of injecting safety guidelines or ethical principles at inference time, allowing alignment without parameter updates. Recent work has positioned RAG as a complement to fine-tuning rather than a substitute, and the relative effectiveness of the two approaches on ethical reasoning benchmarks has received less empirical attention.

\paragraph{Ethical benchmarks.} The ETHICS dataset \citep{hendrycks2021ethics} provides a standardized evaluation across five moral foundations: commonsense morality, justice, deontology, virtue ethics, and utilitarianism. It has become a standard benchmark for assessing whether language models track human ethical judgments.

\paragraph{Small-model alignment.} Most alignment results are reported for models above 7B parameters. The behavior of alignment methods at smaller scale—1--3B parameters—is less well characterized, and prior work has noted that some methods may not transfer cleanly to smaller capacities \citep{xu2024dpo_failures}. Our work targets this regime directly.

\section{Methods}

\subsection{Direct Preference Optimization}
DPO \citep{rafailov2023dpo} minimizes the loss
\begin{equation*}
\mathcal{L}_{\text{DPO}} = -\mathbb{E} \log \sigma \left( \beta \log \frac{\pi_\theta(y_w|x)}{\pi_{\text{ref}}(y_w|x)} - \beta \log \frac{\pi_\theta(y_l|x)}{\pi_{\text{ref}}(y_l|x)} \right)
\end{equation*}
where $(x, y_w, y_l)$ is a preference triple of prompt, preferred completion, and rejected completion; $\pi_\theta$ is the trainable policy; $\pi_{\text{ref}}$ is a frozen reference policy; and $\beta$ controls the strength of the KL constraint to the reference. We use $\beta = 0.1$, in line with the original work.

\subsection{RLHF Pipeline}
Our RLHF implementation follows the canonical pipeline:
\textbf{(1) Reward Model.} We attach a scalar regression head to a Qwen-2.5 backbone and train it via the Bradley-Terry pairwise loss on the same preference data used for DPO, ensuring direct comparability.
\textbf{(2) PPO.} We optimize a policy using PPO \citep{schulman2017ppo} with a KL penalty against the reference policy. The reward signal is the scalar output of the trained reward model evaluated on policy rollouts.

Both stages use LoRA \citep{hu2021lora} adapters with rank $r=16$ targeting all attention projection matrices.

\subsection{Retrieval-Augmented Generation}
We construct a corpus of 21 hand-curated ethical principles spanning major frameworks (utilitarian, deontological, virtue-ethics, justice, common-sense morality, AI safety guidelines, and human rights). Documents are embedded with \texttt{BAAI/bge-small-en} \citep{xiao2024bge} and indexed in FAISS \citep{johnson2017faiss}. At inference, the top-$k$ retrieved principles (we use $k=3$) are prepended to the evaluation prompt before the model produces its judgment.

RAG is composable with both DPO and RLHF: we evaluate \emph{DPO+RAG} and \emph{RLHF+RAG} variants where the fine-tuned policy receives retrieval-augmented prompts.

\section{Experimental Setup}

\subsection{Base Model and Training Data}
\paragraph{Model.} All experiments use \texttt{Qwen2.5-1.5B-Instruct} \citep{qwen2024qwen2} as the base, fine-tuned with LoRA adapters in bf16 precision (PPO uses 4-bit quantization due to memory constraints, see Section~\ref{sec:limitations}).

\paragraph{Preference data.} Our training mixture combines:
\begin{itemize}
    \item \textbf{PKU-SafeRLHF} \citep{ji2024pku} (15{,}000 pairs): safety-oriented preferences from human raters.
    \item \textbf{UltraFeedback} \citep{cui2023ultrafeedback} (5{,}000 pairs): general-quality preferences across multiple dimensions.
\end{itemize}
After length filtering, $17{,}332$ pairs remain (PKU: $14{,}905$; UltraFeedback: $2{,}427$). We do \emph{not} use ETHICS for training; it is held out entirely as the evaluation set.

\paragraph{Hyperparameters.} For all three trainable methods we use LoRA with $r=16$, $\alpha=32$, dropout $0.05$ on attention projection matrices. DPO trains for 1 epoch (1029 steps) with learning rate $5\mathrm{e}{-}5$, $\beta=0.1$. The reward model trains for 1 epoch with learning rate $2\mathrm{e}{-}5$. PPO runs for 1000 episodes with learning rate $1\mathrm{e}{-}5$ and KL coefficient $0.05$, max prompt length $128$, response length $64$.

\subsection{Evaluation Protocol}
We evaluate on the ETHICS test split across all five subsets, using 100 examples per subset (500 total). For classification subsets (commonsense, justice, deontology, virtue), we prompt the model with the canonical formulation for each subset and parse a binary judgment from the first non-whitespace digit in the generation. For utilitarianism, which is pairwise, we score both scenarios under each model and predict the more pleasant one based on average log-likelihood.

We compare six variants:
\textbf{Baseline} (unmodified Qwen2.5-1.5B-Instruct),
\textbf{DPO}, \textbf{RLHF}, \textbf{RAG-only}, \textbf{DPO+RAG}, \textbf{RLHF+RAG}.

\section{Results}

\subsection{Main Results}
Table~\ref{tab:main} reports accuracy on each ETHICS subset and the overall mean across subsets.

\begin{table*}[t]
\centering
\small
\begin{tabular}{l|ccccc|c}
\toprule
\textbf{Model} & \textbf{Commonsense} & \textbf{Justice} & \textbf{Deontology} & \textbf{Virtue} & \textbf{Utilitarianism} & \textbf{Overall} \\
\midrule
Baseline    & 0.510 & 0.500 & 0.540 & 0.400 & 0.480 & 0.486 \\
DPO         & 0.560 & 0.500 & 0.540 & 0.660 & 0.470 & 0.546 \\
RLHF        & 0.520 & 0.500 & 0.540 & 0.610 & 0.470 & 0.528 \\
\midrule
RAG-only    & 0.620 & 0.490 & 0.540 & 0.780 & 0.480 & 0.582 \\
DPO+RAG     & 0.610 & 0.490 & 0.540 & 0.790 & 0.470 & 0.580 \\
RLHF+RAG    & \textbf{0.620} & 0.490 & 0.540 & \textbf{0.790} & 0.470 & \textbf{0.582} \\
\bottomrule
\end{tabular}
\caption{Accuracy on the ETHICS benchmark across five subsets and overall (mean accuracy across subsets, weighted equally). All numbers reported on 100 held-out examples per subset.}
\label{tab:main}
\end{table*}

We highlight three findings:

\paragraph{Finding 1: Retrieval beats preference learning at small scale.}
RAG-only achieves an overall accuracy of $0.582$, exceeding both DPO ($0.546$) and RLHF ($0.528$). Despite involving no parameter updates and using only 21 hand-curated principles, the retrieval baseline outperforms 6.5 hours of DPO training and the full RLHF pipeline. This is consistent with the hypothesis that small models lack the capacity to internalize ethical signals from preference data alone and benefit substantially from explicit context.

\paragraph{Finding 2: DPO outperforms RLHF.}
DPO improves over the baseline by $+6.0$ points; RLHF by only $+4.2$. This ordering matches the recent literature on small-model preference learning, where RLHF's two-stage pipeline introduces compounding noise from a noisy reward model and a sparse RL signal, while DPO's direct optimization is more sample-efficient \citep{rafailov2023dpo,xu2024dpo_failures}.

\paragraph{Finding 3: Composition does not stack.}
Composing retrieval with preference fine-tuning yields no additive gain: DPO+RAG ($0.580$) and RLHF+RAG ($0.582$) match RAG-only ($0.582$) within evaluation noise. This plateau effect suggests that, on ETHICS, retrieval already provides much of the signal that preference fine-tuning would supply, making the two paradigms substantially redundant rather than complementary.

\subsection{Per-Subset Analysis}
The five ETHICS subsets respond very differently to our interventions:
\begin{itemize}
    \item \textbf{Virtue} is the most malleable: baseline accuracy ($0.400$) sits below chance, suggesting systematic confusion about the labeling scheme; RAG-based methods recover to $0.78$--$0.79$, a $+39$-point absolute improvement.
    \item \textbf{Commonsense} responds moderately to all methods, with RAG variants reaching $0.62$.
    \item \textbf{Justice}, \textbf{deontology}, and \textbf{utilitarianism} are essentially flat across all six variants, hovering around the chance-level $0.50$. These subsets appear to demand abstract reasoning capacity that neither retrieval nor preference fine-tuning at this model scale can supply.
\end{itemize}

\subsection{Cost-Benefit Comparison}
RAG requires no training: corpus curation took approximately two hours and FAISS indexing five minutes. DPO required 6.5 hours of GPU time on consumer hardware (RTX 4060 8GB) at bf16 precision. RLHF required an additional 5 hours for the reward model plus 1.5 hours for PPO on cloud GPUs (T4×2). Per absolute point gained on ETHICS, RAG is two orders of magnitude cheaper than the preference-based methods, while delivering superior accuracy.

\section{Discussion}

Our most striking result—retrieval surpassing both DPO and RLHF on ETHICS—has, we believe, three plausible explanations.

\paragraph{Capacity bottleneck.} At 1.5B parameters, the model may lack the representational capacity to encode nuanced ethical distinctions through gradient updates alone. Retrieval bypasses this bottleneck by externalizing the moral framework: rather than asking the model to ``know'' that lying for personal gain is wrong, we provide the relevant principle in-context and ask the model only to apply it.

\paragraph{Distributional mismatch.} Our preference data (PKU + UltraFeedback) is drawn from dialogue scenarios and instruction-following, whereas ETHICS consists of short moral-judgment classifications. Preference fine-tuning teaches the model to produce safer dialogue continuations, which transfers only partially to classifying isolated moral propositions. RAG faces no such mismatch: the retrieved principles are by construction relevant to the evaluation format.

\paragraph{Plateau as redundancy.} The fact that DPO+RAG does not exceed RAG-only suggests that whatever ethical signal preference fine-tuning installs in the policy is already provided more cheaply by retrieval. Composing the two approaches yields diminishing returns, possibly because both ultimately push the model toward the same set of moral conclusions.

\paragraph{Implications.} For practitioners aligning small models under resource constraints, our results suggest that investing in a well-curated retrieval corpus may be more cost-effective than running fine-tuning pipelines. The plateau effect further indicates that hybrid approaches may not always pay for themselves: an effective retrieval system can saturate the alignment improvement available at this scale.

\section{Limitations}
\label{sec:limitations}

We acknowledge several limitations.

\paragraph{Single seed.} All numbers reflect a single random seed; we do not report confidence intervals. The variance of these methods at small scale is non-trivial, and we estimate roughly $\pm 2$--$3$ accuracy points of noise per condition. A multi-seed replication would strengthen the conclusions.

\paragraph{PPO budget and precision.} Compute constraints forced us to run PPO with 4-bit quantization on all four PPO models (policy, reference, value, reward) and to limit training to 1000 episodes. While the QLoRA result \citep{dettmers2023qlora} suggests 4-bit fine-tuning closely matches full precision, the cumulative effect of quantization across four models in a PPO pipeline introduces an additional source of noise that we estimate at $2$--$5$ accuracy points. A bf16 PPO run with a larger episode budget could narrow or close the gap with DPO.

\paragraph{Retrieval corpus scale.} Our hand-curated corpus comprises only 21 documents. A larger, automatically constructed corpus (e.g., 200 or more documents drawn from Wikipedia ethics articles, constitutional AI principles, and the Universal Declaration of Human Rights) could plausibly increase RAG performance further, especially on the subsets where retrieval already dominates.

\paragraph{Single model scale.} We test only the 1.5B parameter regime. Whether the plateau effect and the RAG advantage persist at 3B, 7B, or larger remains an open question. It is plausible that preference fine-tuning becomes more competitive as capacity grows.

\paragraph{No synthetic data.} We did not augment preferences with synthetic ETHICS-format examples generated by a stronger model. Doing so would likely close some of the format gap between PKU/UltraFeedback (dialogue) and ETHICS (short moral classification), but requires either paid API access or significant engineering. We document this as a deliberate omission rather than an oversight.

\section{Ethics Statement}
This work studies methods for aligning language models with widely shared ethical principles. We use only publicly available preference data (PKU-SafeRLHF, UltraFeedback) and a publicly available evaluation benchmark (ETHICS). The retrieval corpus consists of principles drawn from public ethical frameworks and human rights instruments. No personal data was collected or used.

The methods we evaluate are dual-use: the same techniques that align a model with constructive ethical principles could in principle be used to align it with harmful ones. Our experiments report only constructive alignment outcomes; we do not produce or release models trained on harmful preference data.

\section{Conclusion}

We compared three families of ethical alignment methods on a 1.5B-parameter language model: DPO, RLHF (reward model + PPO), and RAG. Across the ETHICS benchmark, RAG—despite using no parameter updates and a hand-curated 21-document corpus—outperformed both preference-based fine-tuning methods, with DPO in turn outperforming RLHF. Composing retrieval with preference fine-tuning produced no additive gain, suggesting the two paradigms encode largely overlapping alignment signals at this scale. We identify five concrete directions for follow-up work: (1) multi-seed replication for confidence intervals, (2) augmenting preferences with synthetic ETHICS-format data, (3) scaling the RAG corpus by an order of magnitude, (4) running PPO at full bf16 precision with a larger episode budget, and (5) testing on ETHICS-hard adversarial split.

% ===== Bibliography =====
\bibliography{references}
\bibliographystyle{acl_natbib}

\end{document}
